\documentclass{article}

\begin{document}
    \begin{center}
        \Large \textbf{Challenge 3: Arbeiten wie ein Power-User} \\ [6pt]
        \normalsize
        Autor: Maximilian Jakob Maag B.Sc. \\
        Version: 1.0
    \end{center}

    \section{Intro}
    Nachdem Sie die Grundlagen kennengelernt haben, geht es nun darum, effizient mit Linux zu arbeiten.
    Der Fokus liegt auf der Nutzung der Shell, Prozessen und einfachen Automatisierungen. \\\\
    In dieser Challenge vertiefen Sie Ihr Wissen im Umgang mit dem System. \\\\

    \section{Lernziele}
    Sie lernen Prozesse zu verwalten und einfache Aufgaben zu automatisieren.
    Verschaffen Sie sich einen Überblick über grundlegende Systemdienste und Netzwerkbefehle.

    \section{Aufgaben}
    Diese Aufgaben sollen Ihnen dabei helfen, den Umgang mit der Shell zu vertiefen und erste Automatisierungen umzusetzen.

    \subsection{Aufgabe 1}
    Was sind Prozesse? Welche Befehle gibt es, um laufende Prozesse anzuzeigen?

    \subsection{Aufgabe 2}
    Starten Sie ein Programm und beenden Sie dieses gezielt über das Terminal.

    \subsection{Aufgabe 3}
    Erstellen Sie ein einfaches Bash-Skript, das mehrere Befehle automatisch ausführt.

    \subsection{Aufgabe 4}
    Nutzen Sie Netzwerkbefehle wie ping oder curl und beschreiben Sie deren Zweck.

    \subsection{Aufgabe 5}
    Informieren Sie sich über Systemdienste und finden Sie heraus, wie man diese startet und stoppt.

\end{document}